\documentclass{article}

\usepackage[utf8]{inputenc}
\usepackage{amssymb}
\usepackage{amsmath}
\usepackage{mathabx}
\usepackage{dcolumn}
\usepackage{geometry}
\usepackage{breqn}
\usepackage{graphicx}
\usepackage{float}
\usepackage{mathrsfs}
\usepackage{array}
\usepackage{caption}
\usepackage{subcaption}
\usepackage[spanish,es-lcroman]{babel}
\decimalpoint
\usepackage{enumerate}
\usepackage{nicefrac} 
\usepackage[most]{tcolorbox}
\usepackage{tcolorbox} % For solution boxes
%\decimalpoint
\setlength\parindent{1.5em}
\usepackage{enumitem}
\newcommand*{\QED}{\hfill\ensuremath{\blacksquare}}
%\usepackage{authblk}
\selectlanguage{spanish}
\geometry{letterpaper, margin=1in}
\pagestyle{headings}
\usepackage{amsthm} 
\newtheorem{theorem}{Teorema}[section]
\newtheorem{corollary}{Corolario}[theorem]
\newtheorem{lemma}[theorem]{Lema}
\theoremstyle{remark}
\newtheorem*{remark}{Considere}
\DeclareUnicodeCharacter{2212}{-}
\usepackage{epigraph}
\usepackage{xcolor}

\usepackage[natbibapa]{apacite}
\usepackage{tikz}
\usepackage{hyperref}

\tcbset{colback=blue!10, 
    %colframe=green!45!black!20, 
    colframe=blue!70!black!70, 
    title=Soluci\'on,
    fonttitle=\bfseries,  
    coltitle=white,
    standard jigsaw, opacityback=50, arc=0mm, breakable} 

\newcommand{\ubar}[1]{\text{\b{$#1$}}}
 
\theoremstyle{definition}
\newtheorem{definition}{Definici\'on}[section]

\title{\textbf{Propuesta Modelo}}

\author{Jhoan Sebasti\'an Fuentes Hern\'andez \and Nicolas Lozano Huertas}
\date{Mayo, 2025}

\begin{document}
\maketitle

Construido sobre Acemoglu \& Restrepo (2022) y Baek \& Jeong (2023)

\section{Introducci\'on}
Se tienen trabajadores con 4 tipos de habilidades: cognitivas, socioemocionales, rutinarias y de juicio, cada una denotada por su inicial. \\
De esta forma definimos el conjunto de tipos de trabajadores: $G=\{j,s,c,r\}$.\\
Cada firma debe decidir si una tarea $x$ la realiza con capital o con trabajo humano.
\section{Firmas}

Las firmas utilizan capital y trabajo para ``producir'' tareas, que posteriormente se utilizan para la producci\'on del bien de cada firma.

\begin{itemize}
    \item \textbf{Producci\'on de tareas}
    \begin{align}
        t(x) =A_{k}\psi_{k}K(x) +A_{r}\psi_{r}l_{r}(x)+A_{s}\psi_{s}l_{s}(x)+A_{j}\psi_{j}l_{j}(x)+A_{c}\psi_{c}l_{c}(x)
    \end{align}
    Donde, $A_{f}$ es la productividad general del factor $f$; $\psi_{f}(x)$ es la productividad espec\'ifica del factor $f$ para producir la tarea $x$; $l_{s}(x), \ l_{c}(x), \ l_{j}(x), \ l_{r}(x), \ K(x)$ son las cantidades de cada tipo de trabajador y de capital utilizadas para producit la tarea x.
    \item \textbf{Funci\'on de producci\'on de la firma:} \\
    Se combina una masa $M_i$ del conjunto de todas las tareas $T$ para producir su bien final. Este mercado funciona bajo competencia monopol\'istica.
    \begin{align*}
        y = A \left[ \frac{1}{M} \int_{T} \left( M \cdot t(x) \right) ^ {\frac{\lambda -1}{\lambda}} dx \right] ^ {\frac{\lambda}{\lambda -1}}
    \end{align*}
    $\lambda$ es la elasticidad de sustituci\'on entre tareas.
    \item \textbf{Costo del capital} \\
    Asumimos que el capital se produce utilizando el bien final con us costo marginal de $\frac{1}{q(x)}$.
    \item \textbf{Regla de asignaci\'on} \\
    Asumimos que para una tarea $x$ siempre se cumple alguna de estas condiciones:
    \begin{align*}
        \frac{w_{g}}{A_{g}\psi_{g}(x)} &< \frac{w_{g'}}{A_{g'}\psi_{g'}(x)} \leq \frac{1}{A_{k}\psi_{k}(x)q(x)} \ ; \forall g \neq g'; \ g,g' \in G \\
        \frac{1}{A_{k}\psi_{k}(x)q(x)} &< \frac{w_{g}}{A_{g}\psi_{g}(x)}; \ g\in G \\
    \end{align*}
    De esta forma garantizamos que un factor productivo siempre es m\'as eficiente que el resto para la realizaci\'on de una tarea, lo que asegura que cada tarea es realizada por un solo factor (el de menor costo relativo). Lo que convierte la ecuaci\'on 1 en:
    \begin{align}
        t(x) =A_{f}\psi_{f}f(x)
    \end{align}
    Donde $f$ es el factor m\'as eficiente para la realizaci\'on de $x$. \\

    Ahora podemos definir los conjuntos $\tau_{f}; \ f \in \{l_{c},l_{r},l_{s},l_{j},K\}$, los cuales contienen las tareas asignadas a cada tipo de trabajo y al capital:

    \begin{align*}
        \tau_{g}&=\left\{x \middle|  \frac{w_{g}}{A_{g}\psi_{g}(x)} < \frac{w_{g'}}{A_{g'}\psi_{g'}(x)} \leq \frac{1}{A_{k}\psi_{k}(x)q(x)} \right\} \\
        \tau_{k}&=\left\{x \middle| \frac{1}{A_{k}\psi_{k}(x)q(x)} < \frac{w_{g}}{A_{g}\psi_{g}(x)} \right\}
    \end{align*}
    \item \textbf{Factor shares} \\
    El share para cada tipo de trabajo est\'a dado por:
    \begin{align*}
        \Gamma_{g} = \frac{1}{M}\int_{{\tau}_{g}}\left(\psi_g(x)\right)^{\lambda-1}dx
    \end{align*}
    El share para el capital esta dado por:
    \begin{align*}
        \Gamma_{k} = \frac{1}{M}\int_{{\tau}_{k}}\left(\psi_k(x)q(x)\right)^{\lambda-1}dx
    \end{align*}
\end{itemize}


\end{document}